%% ma: L12-B08-0189
%% lop: 12
%% chuong: 2
%% bai: 8
%% dang: 12.8.54
%% loai: TN4
%% muc: TH
%% dap_an: "A"
%% nguon: "CĐ2.6.pdf giáo viên gửi 27/10/2026 (M07, Lê Minh Khôi – Nguyễn Anh Khôi, Chương II Chủ đề 6 Ứng dụng vào các bài toán thực tế), Câu 19"
%% kien_thuc: "Bài 8 (tọa độ vectơ, độ dài, tích vectơ với một số)"
%% trang_thai: chinh-thuc
%% ghi_chu: "Giáo viên đã duyệt 27/10/2026"
\begin{ex}
Hai lực $\vec F_1$ và $\vec F_2$ được biểu diễn bởi các vectơ có điểm đầu tại gốc tọa độ. Lực thứ nhất có độ lớn $20$ lb và có hướng đi qua điểm $P(1;1;0)$. Lực thứ hai có độ lớn $40$ lb và có hướng đi qua điểm $Q(0;1;1)$. Gọi $\vec F$ là hợp lực của $\vec F_1$ và $\vec F_2$. Độ lớn của hợp lực $\vec F$ là
\choice{\True $20\sqrt7$ lb}{$60$ lb}{$40\sqrt3$ lb}{$50$ lb}
\loigiai{$\vec F_1=\dfrac{20}{\sqrt2}(1;1;0)=10\sqrt2(1;1;0)$, $\vec F_2=20\sqrt2(0;1;1)$.

$\vec F=10\sqrt2(1;3;2)$, $|\vec F|=10\sqrt2\cdot\sqrt{14}=20\sqrt7$ lb.}
\end{ex}
